%------------------------------------------------------------------------------
% 4.7 تصميم واجهة المستخدم وتجربتها
% المالك: FE (مهمة جانت ٢.٢.٣ — التسليم 04/11/2026)
% المصدر المدموج (2026-11-08):
%   (1) مسودة محمد الشريف (Word «بيان4-1.docx»): مقدمة القسم ونص 4.7.1
%       كاملًا بأنماطه الخمسة، ومسودتا الواجهتين (شاشة المحادثة ولوحة
%       الإدارة) اللتان اشتُقت منهما المخططات السلكية أدناه.
% توحيدات مطبقة على نص المصدر: توحيد ترقيم الأنماط (النمط الأول حتى
%   الخامس) بدل الترقيم المختلط، وإحلال مسميات الوكلاء المعتمدة في 4.2
%   (وكيل SQL، ووكيل RAG، والمُعِدّ) محل الصياغات الحرفية، وتحويل صيغة
%   الاقتراح إلى قرار تصميم مسجل (زر تصدير التقرير)، وتغليف المصطلحات
%   اللاتينية بـ \EN{}.
% وفي المخططين السلكيين صُوّبت لغويًا تسميات وردت في مسودة الواجهات
%   («قبول» بدل «قبون»، و«الأدوار» بدل «الأبار»، و«إحصائيات» بدل
%   «إجصاصات»)، واستُكمل الناقص فيها: لوحة «العروض الذكية»
%   الغائبة من لوحة الإدارة، ومؤشر الانتظار في شاشة المحادثة، وحقل
%   الإدخال وزر الإرسال المقتطعان منها.
% التتبع: يخدم حالات الاستخدام UC-01 حتى UC-06، والمتطلبات FR-02 وFR-05
%   وFR-06 وFR-12 وFR-18 وFR-23 وFR-24 وFR-26 وFR-29، وغير الوظيفية
%   NFR-13 وNFR-14 وNFR-15؛ وتُجمع الروابط في مصفوفة القسم 4.8.
% ملاحظة: التصميم المرئي (الهوية والألوان والخطوط) موضعه الطبيعي الفصل
%   الخامس عند عرض التنفيذ — هنا مستوى الهيكل وتدفق المستخدم فقط.
% حالة الإنجاز: مكتمل (نص 4.7.1 مدموج + مخططان سلكيان)؛ ينتظر المراجعة
%   المتقاطعة والتجميع النهائي.
%------------------------------------------------------------------------------
\section{تصميم واجهة المستخدم وتجربتها}
\label{sec:ui-design}

يركّز تصميم واجهة المستخدم في نظام «بيان» على توفير تجربة استخدام عربية واضحة وسلسة تمكّن المستخدم من التفاعل مع النظام باللغة العربية بطريقة طبيعية، مع مراعاة اتجاه العرض من اليمين إلى اليسار (\EN{RTL})، ووضوح مخرجات النظام، وإبراز عناصر الشفافية مثل الاستعلامات المنفذة والمصادر المرجعية عند الحاجة. كما يهدف التصميم إلى دعم كل من المستخدم العادي والمستخدم المتقدم من خلال واجهات مبسطة في الاستخدام، مع المحافظة على عناصر التحكم الإداري والأمني اللازمة لإدارة المستخدمين، وتتبع العمليات، ومراقبة أداء النظام.

ويتوزع هذا القسم على فقرتين: تدفق الحوار العربي بأنماطه الخمسة المعتمدة (الفقرة \ref{subsec:ui-conversation})، والمخططات السلكية للشاشتين الرئيستين (الفقرة \ref{subsec:ui-wireframes}). ويخدم التصميم مباشرةً حالات الاستخدام \EN{UC-01} حتى \EN{UC-06} من الفصل الثالث، والمتطلبات الوظيفية \EN{FR-02} و\EN{FR-05} و\EN{FR-06} و\EN{FR-12} و\EN{FR-18} و\EN{FR-23} و\EN{FR-24} و\EN{FR-26} و\EN{FR-29}، وغير الوظيفية \EN{NFR-13} و\EN{NFR-14} و\EN{NFR-15}؛ وتُجمع روابطه في مصفوفة التتبع بالقسم \ref{sec:design-traceability}. أما الهوية البصرية — الألوان والخطوط والإيقاع — فموضعها الطبيعي فصل التنفيذ والاختبار، إذ يقف هذا القسم عند مستوى الهيكل وتدفق المستخدم.

\subsection{تدفق الحوار العربي وأنماط المحادثة}
\label{subsec:ui-conversation}

يعتمد تدفق الحوار على خمسة أنماط معتمدة تحكم كل جلسة محادثة، ويجسدها الشكل \ref{fig:wireframes-chat} بأرقامها من (1) إلى (5):

\textbf{النمط الأول — اقتراحات الأسئلة الافتتاحية عند بدء المحادثة:}\\
عند بدء المحادثة في نظام «بيان»، تعرض الواجهة مجموعة من اقتراحات الأسئلة الافتتاحية لمساعدة المستخدم على البدء بسرعة ووضوح، خاصة إذا كان يستخدم النظام لأول مرة. تظهر هذه الاقتراحات بصياغات عربية مباشرة ومناسبة لطبيعة العمل المؤسسي، مثل: «كم عدد الطلاب المسجلين في كلية الهندسة؟»، أو «اعرض لائحة الغياب والحرمان»، أو «أنشئ تقريرًا عن أداء الأقسام خلال الفصل الحالي». ويساهم هذا الأسلوب في تقليل التردد عند بداية الاستخدام، كما يعرّف المستخدم ضمنيًا بنطاق قدرات النظام، سواء في الاستعلام عن البيانات أو البحث في اللوائح أو توليد التقارير. وقد روعي في عرض هذه الاقتراحات أن تكون مختصرة وواضحة ومتوافقة مع الواجهة العربية ذات الاتجاه من اليمين إلى اليسار، بما يعزز سهولة الاستخدام وسرعة التعلم.

\textbf{النمط الثاني — عرض استعلام \EN{SQL} المنفذ للمستخدم المتقدم:}\\
يراعي تصميم واجهة المحادثة في نظام «بيان» مبدأ الشفافية التقنية من خلال إتاحة الاطلاع على استعلام \EN{SQL} المنفذ للمستخدمين المتقدمين المصرّح لهم، وذلك بهدف تعزيز موثوقية النتائج وتسهيل التحقق من آلية استخراج البيانات. فعند طرح المستخدم سؤالًا باللغة العربية يتعلق بالبيانات المؤسسية، يتولى وكيل \EN{SQL} تحويل السؤال إلى استعلام مناسب بالاعتماد على طبقة العروض الذكية (\EN{Smart Views})، ثم تنفيذه ضمن حدود الصلاحيات الأمنية المحددة. وتُعرض النتيجة للمستخدم بصورة مبسطة، مع توفير خيارٍ قابلٍ للتوسيع بعنوان «عرض استعلام \EN{SQL}» يُظهر نص الاستعلام المنفذ دون التأثير في سير المحادثة. ويكون هذا العرض مخصصًا للاطلاع والتحقق فقط، مع مراعاة التحكم في الوصول المبني على الأدوار (\EN{RBAC})، بما يضمن عدم كشف تفاصيل البيانات أو الاستعلامات للمستخدمين غير المصرّح لهم. ويسهم هذا الأسلوب في تحقيق التوازن بين سهولة الاستخدام للمستخدم العادي والشفافية التقنية للمستخدم المتقدم، مع دعم قابلية التدقيق وتتبع العمليات داخل النظام.

\textbf{النمط الثالث — إظهار المصادر وأرقام الصفحات في إجابات الاسترجاع:}\\
يعتمد تصميم واجهة المحادثة في نظام «بيان» على مبدأ توثيق المعلومات المسترجعة من المستندات واللوائح المؤسسية، بما يعزز موثوقية الإجابات وقابليتها للتحقق. فعندما يطرح المستخدم سؤالًا يتعلق بالسياسات أو اللوائح التنظيمية، يتولى وكيل الاسترجاع المعزز بالتوليد (\EN{RAG}) البحث الدلالي في المستندات المفهرسة باستخدام قاعدة البيانات المتجهية (\EN{Qdrant})، واسترجاع المقاطع النصية ذات الصلة بالسؤال. وتُعرض الإجابة الناتجة باللغة العربية مع إظهار المراجع المرتبطة بها، بما يشمل اسم المستند ورقم الصفحة التي استُخلصت منها المعلومات، وذلك في مساحة واضحة أسفل الإجابة أو بجانبها. ويتيح هذا التصميم للمستخدم تتبع مصدر المعلومة والرجوع إلى الوثيقة الأصلية عند الحاجة، مع مراعاة صلاحيات الوصول إلى المستندات. وفي حال عدم توفر مصادر كافية تدعم الإجابة، تعرض الواجهة رسالة توضيحية مناسبة بدلًا من تقديم معلومات غير موثقة. ويسهم هذا الأسلوب في الحد من مخاطر الإجابات غير الدقيقة، وتعزيز الشفافية والثقة في استخدام المساعد الذكي داخل البيئة المؤسسية.

\textbf{النمط الرابع — تصدير التقرير من نتيجة الاستعلام:}\\
يراعي تصميم واجهة المستخدم في نظام «بيان» توفير آلية واضحة ومباشرة لتصدير التقارير الناتجة عن الاستعلامات، بما يتيح للمستخدم الاستفادة من البيانات المستخرجة خارج بيئة المحادثة. فعند طرح سؤال يتطلب استخراج بيانات أو إعداد تقرير تحليلي، يتولى الوكيل المختص معالجة الاستعلام واسترجاع النتائج، بينما يعمل المُعِدّ (\EN{Reporter}) على تنظيم المخرجات وتنسيقها في صورة مناسبة للعرض، سواء أكانت نصوصًا أم جداول أم رسومًا بيانية بحسب طبيعة النتائج. وتتضمن واجهة المحادثة زرًا بعنوان «تصدير التقرير» يظهر عند توفر نتائج قابلة للتصدير، ويتيح للمستخدم اختيار صيغة الملف المناسبة، مثل \EN{PDF} أو \EN{Excel} أو \EN{CSV}. وتُراعى في هذه العملية صلاحيات المستخدم وفق آلية التحكم في الوصول المبني على الأدوار (\EN{RBAC})، بما يضمن عدم تصدير بيانات تتجاوز نطاق الصلاحيات الممنوحة له. ويسهم هذا التصميم في تسهيل إعداد التقارير المؤسسية، وتقليل الحاجة إلى إعادة إدخال البيانات أو معالجتها يدويًا، إضافةً إلى دعم عمليات التوثيق والمتابعة واتخاذ القرارات بالاعتماد على معلومات منظمة وقابلة للاستخدام.

\textbf{النمط الخامس — حالات الانتظار والأخطاء برسائل عربية واضحة:}\\
يراعي تصميم واجهة المحادثة في نظام «بيان» توفير تغذية راجعة مستمرة للمستخدم أثناء معالجة الطلبات، بما يضمن وضوح حالة النظام ويعزز سهولة التفاعل معه. فعند إرسال المستخدم سؤالًا، تُعرض رسالة انتظار باللغة العربية مثل «جارٍ معالجة طلبك...» مع مؤشر بصري يوضح استمرار العملية، ويمكن إظهار رسائل أكثر تحديدًا بحسب طبيعة المهمة، مثل «جارٍ البحث في المستندات واللوائح...» أو «جارٍ إعداد التقرير...». وعند اكتمال المعالجة، تظهر الإجابة أو النتائج بصورة واضحة داخل المحادثة. أما عند حدوث خطأ، فيُعرض للمستخدم رسالة عربية مبسطة توضح طبيعة المشكلة دون إظهار التفاصيل البرمجية الحساسة، مثل «تعذر تنفيذ الطلب، يُرجى المحاولة مرة أخرى»، مع إتاحة خيار إعادة المحاولة عند الحاجة. كما تميز الواجهة بين أخطاء الاتصال، وعدم كفاية الصلاحيات، وعدم العثور على نتائج مطابقة، بحيث يحصل المستخدم على توجيه مناسب لكل حالة. ويسهم هذا التصميم في تحسين تجربة المستخدم وتقليل الغموض أثناء تنفيذ العمليات، مع تعزيز موثوقية النظام وقابليته للاستخدام في البيئة المؤسسية.

\subsection{المخططات السلكية للشاشات الرئيسة}
\label{subsec:ui-wireframes}

تُثبَّت الشاشتان الرئيستان عند مستوى المخطط السلكي: صناديق رمادية بسيطة بلا هوية بصرية، وفق عرف التوثيق المعتمد في هذا الفصل؛ فالقرارات البصرية تُؤجل إلى فصل التنفيذ والاختبار حيث تُقاس قابلية الاستخدام (\EN{NFR-13}، \EN{NFR-14}). والشاشة الأولى هي واجهة المحادثة الرئيسة التي يمر بها كل طلب (الشكل \ref{fig:wireframes-chat})، وقد اشتُقت من مسودة الواجهات الأولية للفريق بعد تصويب تسمياتها اللغوية واستكمال الناقص فيها: حقل الإدخال وزر الإرسال، ومؤشر الانتظار، وأرقام الأنماط من (1) إلى (5) المربوطة بأنماط الفقرة \ref{subsec:ui-conversation}. والشاشة الثانية هي لوحة الإدارة (الشكل \ref{fig:wireframes-admin}) بمهامها الثلاث كاملة، وهي من نصيب مسؤول تقنية المعلومات بالمؤسسة: فإدارة المستخدمين والأدوار تعرض جدول الحسابات بالدور والحالة وإجراء «تعديل الدور» (\EN{UC-04})؛ وسجل التدقيق يعرض الوقت والمستخدم والإجراء والنتيجة، ومنه تُقرأ حالات الرفض الأمني كإحباط استعلام تجاوز صلاحية صاحبه (\EN{UC-06})؛ ولوحة «العروض الذكية» — التي اقتصرت مسودة الواجهة على عنصر تنقل لها — تعرض بطاقات الكتالوج بتصنيفها الأمني وحالة اعتمادها مع إجراء «إنشاء عرض ذكي»، فتُدار فيها دورة حياة العرض من الإنشاء إلى الاعتماد وفق ما فصله القسم \ref{sec:smart-views-design} (\EN{UC-05}).

\begin{figure}[H]
  \centering
  %==============================================================================
% figures/tikz/wireframes-chat.tex
% Wireframe of the main chat screen — Chapter 4, Section 4.7.2 (fig:wireframes-chat)
%==============================================================================
% Grayscale wireframe of BAYAN's main conversation screen, rebuilt in TikZ
% from the team's original UI draft after the linguistic fixes recorded in
% sec7_ui.tex («قبول» not «قبون», «الأدوار», «إحصائيات») and after
% completing what the draft was missing: the input field + send button, the
% waiting indicator, and the numbered pattern markers (1)-(5) tied to the
% five dialogue patterns of subsection 4.7.1.
% Reading order is right-to-left — the original designer's choice kept
% intact: navigation sidebar at the far right, user bubble on the right,
% assistant messages on the left.
% Style ratio: wireframe greys only — gray!5..25 / kaulight fills,
% gray!60/70 / dividersilver strokes, kaudark text. The only semantic
% exceptions allowed by the house rule: the kauprimary number circles
% (1)-(5), the single green!60!black connection dot, and the dashed
% red!70!black error box. Arabic inside \textarabic{}, Latin inside \EN{},
% SQL inside \CODE{} (monospace, LTR — the ellipses are deliberate
% redactions, no view/table names are invented). Image/logo placeholders
% are boxes with crossed diagonals; no icons, no pictures, no decorations.
%==============================================================================
\resizebox{\textwidth}{!}{%
\begin{tikzpicture}[
  >=stealth,
  font=\scriptsize,
  text=kaudark,
  % ---- screen chrome ------------------------------------------------------
  frame/.style={draw=gray!70, line width=1pt, fill=white, rounded corners=6pt},
  panel/.style={fill=gray!6, rounded corners=6pt},
  strip/.style={fill=gray!6, draw=gray!60, line width=0.6pt, rounded corners=4pt},
  field/.style={fill=white, draw=gray!60, line width=0.6pt, rounded corners=3pt},
  % ---- chat bubbles (assistant = left, user = right) -----------------------
  bubbleA/.style={fill=gray!10, draw=gray!60, line width=0.6pt, rounded corners=6pt,
                  inner sep=6pt, align=right},
  bubbleU/.style={fill=gray!20, draw=gray!60, line width=0.6pt, rounded corners=6pt,
                  inner sep=6pt, align=right},
  sugcard/.style={fill=white, draw=gray!60, line width=0.6pt, rounded corners=3pt,
                  inner sep=3.5pt, text width=6.6cm, align=right},
  % ---- sidebar navigation --------------------------------------------------
  navitem/.style={fill=white, draw=gray!60, line width=0.6pt, rounded corners=3pt,
                  inner xsep=3pt, inner ysep=3pt, text width=2.45cm, align=right},
  navactive/.style={navitem, fill=gray!20, draw=gray!70, font=\scriptsize\bfseries},
  % ---- wireframe buttons ---------------------------------------------------
  wbtn/.style={fill=white, draw=gray!60, line width=0.6pt, rounded corners=3pt,
               inner xsep=5pt, inner ysep=3pt},
  gbtn/.style={wbtn, fill=gray!20},
  % ---- pattern number circles (semantic exception: kauprimary) ------------
  num/.style={circle, draw=kauprimary, fill=kauprimary!8, text=kaudark,
              font=\scriptsize\bfseries, inner sep=1pt, minimum size=0.42cm},
]

%==============================================================================
% SCREEN FRAME + RIGHT SIDEBAR (RTL: navigation lives at the far right)
%==============================================================================
\path[frame] (0.4,0.6) rectangle (15.6,13.6);
\path[panel] (12.5,0.6) rectangle (15.5,13.6);

% --- sidebar header: logo placeholder (crossed-diagonal box) + branding
\draw[field] (13.3,12.45) rectangle (14.7,13.35);
\draw[draw=gray!60, line width=0.5pt]
  (13.3,12.45) -- (14.7,13.35) (13.3,13.35) -- (14.7,12.45);
\node[font=\small\bfseries] at (14.0,12.2) {\textarabic{بيان}};
\node at (14.0,11.85) {\textarabic{مساعد المؤسسة}};

% --- navigation items (right-aligned; active item filled + bold)
\node[navactive, anchor=east] at (15.32,11.3) {\textarabic{الرئيسية}};
\node[anchor=east, text=gray] at (15.3,10.68) {\textarabic{المحادثات السابقة}};
\node[navitem, anchor=east] at (15.32,10.15) {\textarabic{بطاقة الطالب — كلية الهندسة}};
\node[navitem, anchor=east] at (15.32,9.15) {\textarabic{لائحة الغياب والحرمان}};
\node[navitem, anchor=east] at (15.32,8.35) {\textarabic{التقارير}};
\node[navitem, anchor=east] at (15.32,7.7) {\textarabic{الإعدادات}};

% --- sidebar footer: user card + logout
\path[fill=gray!12, draw=gray!60, line width=0.6pt, rounded corners=3pt]
  (12.66,0.8) rectangle (15.34,2.0);
\node at (14.0,1.63) {\textarabic{أحمد — مدير المؤسسة}};
\node[wbtn, inner ysep=2.5pt] at (14.0,1.12) {\textarabic{تسجيل خروج}};

%==============================================================================
% TOP BAR: connection status (right) / current-chat title (center) /
%          new-chat button (left)
%==============================================================================
\path[strip] (0.6,12.65) rectangle (12.3,13.55);
\fill[green!60!black] (12.05,13.1) circle (0.055);% the single green status dot
\node[anchor=east] at (11.9,13.1) {\textarabic{النظام متصل}};
\node[font=\small\bfseries] at (6.45,13.1) {\textarabic{المحادثة الحالية}};
\node[gbtn] at (1.85,13.1) {\textarabic{محادثة جديدة}};

%==============================================================================
% (1) WELCOME BUBBLE + OPENING SUGGESTIONS (assistant side = left)
%==============================================================================
\node[bubbleA, text width=6.6cm, anchor=north east] at (7.9,12.5)
  {\textarabic{مرحبًا! أنا بيان، مساعدك المؤسسي. يمكنك البدء بالأسئلة التالية:}};
\node[sugcard, anchor=north east] at (7.9,11.45)
  {\textarabic{كم عدد الطلاب المسجلين في كلية الهندسة؟}};
\node[sugcard, anchor=north east] at (7.9,10.8)
  {\textarabic{اعرض لائحة الغياب والحرمان}};
\node[sugcard, anchor=north east] at (7.9,10.15)
  {\textarabic{أنشئ تقريرًا عن أداء الأقسام خلال الفصل الحالي}};
\node[num] at (8.15,11.05) {\EN{1}};

%==============================================================================
% USER BUBBLE (user side = right)
%==============================================================================
\node[bubbleU, text width=6.4cm, anchor=north east] at (12.2,9.28)
  {\textarabic{اعرض لائحة الغياب والحرمان لكلية الهندسة في الفصل الأول}};

%==============================================================================
% ASSISTANT REPLY BLOCK: answer + (2) SQL disclosure + (3) sources + (4) export
%==============================================================================
\path[fill=gray!10, draw=gray!60, line width=0.6pt, rounded corners=5pt]
  (0.9,4.42) rectangle (9.7,8.4);

% answer sentence
\node[anchor=north east, inner sep=0pt, text width=7.9cm, align=right] at (9.45,8.25)
  {\textarabic{إليك لائحة الغياب والحرمان لكلية الهندسة في الفصل الأول، مبنية على بيانات الفصل الحالي.}};

% (2) collapsible SQL bar + LTR code box (ellipses = deliberate redactions)
\node[field, minimum width=3.9cm, minimum height=0.52cm, anchor=east] (sqlbar) at (9.45,7.17) {};
\node[anchor=east] at (9.27,7.17) {\textarabic{عرض استعلام \EN{SQL}}};
\node[regular polygon, regular polygon sides=3, shape border rotate=-90,
      fill=gray!50, inner sep=0pt, minimum size=0.17cm] at (5.95,7.17) {};
\node[num] at (9.95,7.17) {\EN{2}};
\path[fill=gray!15, draw=gray!60, line width=0.6pt, rounded corners=2pt]
  (1.15,6.15) rectangle (9.45,6.79);
\node[anchor=west] at (1.35,6.58) {\CODE{SELECT … FROM v\_…}};
\node[anchor=west] at (1.35,6.34) {\CODE{WHERE college = \$1 ORDER BY absence\_ratio DESC}};

% (3) sources with page numbers (small gray square markers, RTL)
\node[anchor=east, font=\scriptsize\bfseries] at (9.45,5.85) {\textarabic{المصادر والمراجع}};
\fill[gray!45] (9.28,5.41) rectangle (9.4,5.53);
\node[anchor=east] at (9.18,5.47) {\textarabic{اللائحة الأكاديمية للبكالوريوس — ص 12}};
\fill[gray!45] (9.28,5.05) rectangle (9.4,5.17);
\node[anchor=east] at (9.18,5.11) {\textarabic{دليل قبول الطلاب — ص 8}};
\node[num] at (9.95,5.85) {\EN{3}};

% (4) export row: label + PDF / Excel / CSV buttons
\node[anchor=east] at (9.45,4.75) {\textarabic{تصدير التقرير:}};
\node[wbtn, inner ysep=2pt] at (6.55,4.75) {\EN{PDF}};
\node[wbtn, inner ysep=2pt] at (5.3,4.75) {\EN{Excel}};
\node[wbtn, inner ysep=2pt] at (4.0,4.75) {\EN{CSV}};
\node[num] at (9.95,4.75) {\EN{4}};

%==============================================================================
% (5a) WAITING BUBBLE: message + three-dot indicator + progress hint
%==============================================================================
\path[bubbleA] (0.9,3.40) rectangle (5.2,4.28);
\node[anchor=east] at (5.0,4.07) {\textarabic{جارٍ معالجة طلبك...}};
\fill[gray!50] (4.12,3.86) circle (0.05) (4.37,3.86) circle (0.05) (4.62,3.86) circle (0.05);
\node[anchor=east, text=gray] at (5.0,3.66) {\textarabic{جارٍ البحث في المستندات واللوائح...}};
\node[num] at (5.45,3.84) {\EN{5}};

%==============================================================================
% (5b) ERROR BOX: dashed red frame + drawn warning triangle + retry button
%==============================================================================
\path[draw=red!70!black, dashed, line width=0.8pt, fill=red!4, rounded corners=4pt]
  (0.9,1.95) rectangle (7.6,3.1);
\draw[draw=red!70!black, line width=0.7pt, fill=white]
  (6.98,2.56) -- (7.38,2.56) -- (7.18,2.93) -- cycle;
\node[text=red!70!black, font=\scriptsize\bfseries] at (7.18,2.67) {\EN{!}};
\node[anchor=east] at (6.85,2.75) {\textarabic{تعذر تنفيذ الطلب، يُرجى المحاولة مرة أخرى}};
\node[gbtn, inner ysep=2pt] at (2.05,2.26) {\textarabic{إعادة المحاولة}};
\node[num] at (7.85,2.5) {\EN{5}};

%==============================================================================
% INPUT BAR: text field with faint placeholder (right) + send button (left)
%==============================================================================
\path[strip] (0.6,0.7) rectangle (12.3,1.75);
\path[field] (2.75,0.9) rectangle (12.1,1.55);
\node[anchor=east, text=gray] at (11.9,1.225) {\textarabic{اكتب سؤالك بالعربية...}};
\node[fill=gray!25, draw=gray!60, line width=0.6pt, rounded corners=3pt,
      font=\scriptsize\bfseries, inner xsep=8pt, inner ysep=4pt] at (1.65,1.225)
  {\textarabic{إرسال}};

%==============================================================================
% KEY LINE (below the screen frame): the five numbered dialogue patterns
%==============================================================================
\node[inner sep=1pt] at (8,0.18) {\textarabic{%
  \textbf{\textcolor{kauprimary}{\EN{1}}}~الاقتراحات الافتتاحية\hspace{0.4cm}%
  \textbf{\textcolor{kauprimary}{\EN{2}}}~عرض استعلام \EN{SQL}\hspace{0.4cm}%
  \textbf{\textcolor{kauprimary}{\EN{3}}}~المصادر وأرقام الصفحات\hspace{0.4cm}%
  \textbf{\textcolor{kauprimary}{\EN{4}}}~تصدير التقرير\hspace{0.4cm}%
  \textbf{\textcolor{kauprimary}{\EN{5}}}~الانتظار والأخطاء}};

\end{tikzpicture}%
}

  \caption{المخطط السلكي لشاشة المحادثة الرئيسة بواجهة عربية من اليمين إلى اليسار؛ الأرقام من (1) إلى (5) تشير إلى أنماط الحوار الخمسة المعتمدة في الفقرة \ref{subsec:ui-conversation}: الاقتراحات الافتتاحية، وعرض استعلام \EN{SQL}، والمصادر وأرقام الصفحات، وتصدير التقرير، وحالات الانتظار والأخطاء.}
  \label{fig:wireframes-chat}
\end{figure}

\begin{figure}[H]
  \centering
  %==============================================================================
% figures/tikz/wireframes-admin.tex
% Admin-panel wireframe (grayscale) — Chapter 4, Section 4.7.2
% (fig:wireframes-admin)
%==============================================================================
% Low-fidelity wireframe of the administration console used by the Tenant
% IT Administrator. Right-hand sidebar (RTL-first): logo placeholder,
% five navigation items with «إدارة المستخدمين والأدوار» active, and the
% bottom user card with logout. The content area stacks the three admin
% tasks vertically: (1) users & roles management with the accounts table
% [UC-04], (2) the audit log with date/type filters and the security
% rejection row [UC-06], and (3) the Smart Views console that was absent
% from the original interface draft — catalog cards with technical name,
% brief description, security tier, approval status, and review/approve
% actions [UC-05]. Serves the paragraph between fig:wireframes-chat and
% fig:wireframes-admin in sec7_ui.tex (three admin tasks of the console).
% Spelling fixes over the interface draft: «الأدوار» (not «الأبار») and
% «سجل التدقيق» (not «التدفق»); official role names verbatim from the
% stakeholders table (sec3_requirements.tex); the three view names and
% security tiers verbatim from the Smart Views catalog
% (sec3_smart_views.tex, tab:sv-catalog).
% Style: grayscale-only wireframe per the 4.7.2 convention (صناديق رمادية
% بلا هوية بصرية) — tabular-in-node tables as in bayan-er.tex,
% \textarabic{} for Arabic text, \EN{}/\CODE{} for Latin/monospace, and
% right-to-left reading order (sidebar far right; first table column on
% the right). Semantic exceptions only: the faded \EN{UC-xx} tags and the
% red «مرفوض (تجاوز الصلاحية)» cell in the audit log.
%==============================================================================
\resizebox{\textwidth}{!}{%
\begin{tikzpicture}[
  >=stealth,
  panel/.style={rounded corners=5pt, draw=gray!70, line width=0.9pt},
  field/.style={rounded corners=3pt, draw=gray!60, line width=0.7pt,
    fill=gray!10},
  btn/.style={rounded corners=3pt, draw=gray!70, line width=0.7pt,
    fill=gray!15, font=\scriptsize, text=kaudark, inner sep=2.5pt},
  btnprim/.style={btn, fill=gray!25, font=\scriptsize\bfseries},
  navpill/.style={rounded corners=3pt, draw=gray!60, line width=0.7pt,
    fill=gray!13, inner sep=0pt, minimum width=2.72cm,
    minimum height=0.52cm},
  navtext/.style={anchor=east, font=\scriptsize, text=kaudark,
    inner sep=0pt},
  lblE/.style={anchor=east, font=\scriptsize, text=kaudark,
    inner sep=0pt},
  ttl/.style={anchor=east, font=\small\bfseries, text=kaudark,
    inner sep=0pt},
  uclbl/.style={anchor=west, font=\scriptsize, text=gray!70,
    inner sep=0pt},
  ph/.style={anchor=east, font=\scriptsize, text=gray, inner sep=0pt},
  viewcard/.style={rounded corners=4pt, draw=gray!65, line width=0.8pt,
    fill=gray!8},
  cardtxt/.style={font=\scriptsize, text=kaudark, inner sep=1pt},
  cbtn/.style={rounded corners=3pt, draw=gray!70, line width=0.7pt,
    fill=gray!15, font=\scriptsize, text=kaudark, inner sep=0pt,
    minimum width=1.35cm, minimum height=0.42cm},
]

  %==========================================================================
  % SCREEN FRAME (fill first) + RIGHT SIDEBAR (RTL: sidebar far right)
  %==========================================================================
  \fill[white, rounded corners=8pt] (0.4,0.5) rectangle (15.6,13.1);
  \begin{scope}
    \clip[rounded corners=8pt] (0.4,0.5) rectangle (15.6,13.1);
    \fill[gray!6] (12.5,0.5) rectangle (15.6,13.1);
  \end{scope}
  \draw[gray!60, line width=0.6pt] (12.5,0.5) -- (12.5,13.1);

  %---- Logo placeholder (box with crossed diagonals) + brand ---------------
  \draw[gray!70, line width=0.7pt, fill=white]
    (13.62,12.32) rectangle (14.38,12.94);
  \draw[gray!70, line width=0.5pt]
    (13.62,12.32) -- (14.38,12.94) (13.62,12.94) -- (14.38,12.32);
  \node[font=\small\bfseries, text=kaudark] at (14,12.08)
    {\textarabic{بيان}};
  \node[font=\scriptsize, text=kaudark] at (14,11.80)
    {\textarabic{لوحة الإدارة}};

  %---- Navigation items (active = إدارة المستخدمين والأدوار) ---------------
  \node[navpill] at (14,11.30) {};
  \node[navtext] at (15.18,11.30) {\textarabic{الرئيسية}};

  \node[navpill, fill=gray!20, draw=gray!70, minimum height=0.72cm]
    at (14,10.45) {};
  \node[navtext, font=\scriptsize\bfseries] at (15.18,10.59)
    {\textarabic{إدارة المستخدمين}};
  \node[navtext, font=\scriptsize\bfseries] at (15.18,10.31)
    {\textarabic{والأدوار}};

  \node[navpill] at (14,9.60) {};
  \node[navtext] at (15.18,9.60) {\textarabic{العروض الذكية}};

  \node[navpill] at (14,8.75) {};
  \node[navtext] at (15.18,8.75) {\textarabic{سجل التدقيق}};

  \node[navpill] at (14,7.90) {};
  \node[navtext] at (15.18,7.90) {\textarabic{الإعدادات}};

  %---- Bottom user card (full official role name) + logout ------------------
  \path[rounded corners=3pt, draw=gray!60, line width=0.7pt, fill=gray!12]
    (12.64,0.72) rectangle (15.36,2.05);
  \node[anchor=north, text width=2.45cm, align=center, font=\scriptsize,
        text=kaudark, inner sep=0pt] at (14,1.93)
    {\textarabic{سارة — مسؤول تقنية المعلومات بالمؤسسة}};
  \node[btn, fill=white, draw=gray!60, minimum width=1.6cm,
        minimum height=0.40cm] at (14,1.02) {\textarabic{تسجيل خروج}};

  %==========================================================================
  % PANEL 1 (top): إدارة المستخدمين والأدوار [UC-04]
  %==========================================================================
  \path[panel, fill=gray!5, draw=none] (0.6,8.75) rectangle (12.3,12.95);
  \begin{scope}
    \clip[rounded corners=5pt] (0.6,8.75) rectangle (12.3,12.95);
    \fill[gray!20] (0.6,12.42) rectangle (12.3,12.95);
  \end{scope}
  \draw[panel] (0.6,8.75) rectangle (12.3,12.95);
  \node[ttl] at (12.1,12.68) {\textarabic{إدارة المستخدمين والأدوار}};
  \node[uclbl] at (0.8,12.68) {\EN{UC-04}};

  %---- Toolbar: search field (right) + add-user button (far left) -----------
  \path[field] (7.0,11.85) rectangle (12.1,12.30);
  \node[ph] at (11.95,12.08) {\textarabic{بحث بالاسم أو البريد...}};
  \node[btnprim, minimum width=1.75cm, minimum height=0.45cm]
    at (1.675,12.08) {\textarabic{إضافة مستخدم}};

  %---- Accounts table (tabular-in-node; columns ordered from the right). ----
  % bidi mirrors tabular: the FIRST source column renders rightmost, so the
  % source is written directly in the intended visual right-to-left order
  % (الاسم rightmost ... إجراءات leftmost). Cells are flushed right
  % (\raggedleft) for the RTL table look, and every column is wide enough to
  % keep its longest Arabic string on ONE line (a wrapped multi-line
  % \textarabic paragraph inside an LTR p-column gets its second line
  % displaced — bidi/p-column interplay), so single-line cells are a hard
  % requirement here (widths verified against the 150-dpi render).
  \node[anchor=north east, inner sep=0pt] at (12.1,11.62) {%
    {\footnotesize
    \renewcommand{\arraystretch}{1.3}%
    \begin{tabular}{@{}>{\raggedleft\arraybackslash}p{1.50cm}@{\hspace{6pt}}>{\raggedleft\arraybackslash}p{3.40cm}@{\hspace{6pt}}>{\raggedleft\arraybackslash}p{2.90cm}@{\hspace{6pt}}>{\raggedleft\arraybackslash}p{1.15cm}@{\hspace{6pt}}>{\raggedleft\arraybackslash}p{1.45cm}@{}}
    \arrayrulecolor{gray!50}
    \rowcolor{gray!15}
    \textarabic{\textbf{الاسم}} & \textarabic{\textbf{البريد المؤسسي}} &
    \textarabic{\textbf{الدور}} & \textarabic{\textbf{الحالة}} &
    \textarabic{\textbf{إجراءات}} \\ \hline
    \textarabic{أحمد عبدالله} & \EN{ahmed@example.edu} &
    \textarabic{مدير المؤسسة} & \textarabic{نشط} &
    \textarabic{تعديل الدور} \\ \hline
    \textarabic{سارة محمد} & \EN{sara@example.edu} &
    \textarabic{مسؤول تقنية المعلومات} & \textarabic{نشط} &
    \textarabic{تعديل الدور} \\ \hline
    \textarabic{خالد علي} & \EN{khaled@example.edu} &
    \textarabic{موظف المؤسسة} & {\scriptsize\textarabic{[موقوف]}} &
    \textarabic{تعديل الدور} \\ \hline
    \end{tabular}}};

  %==========================================================================
  % PANEL 2 (middle): سجل التدقيق [UC-06]
  %==========================================================================
  \path[panel, fill=gray!5, draw=none] (0.6,4.85) rectangle (12.3,8.53);
  \begin{scope}
    \clip[rounded corners=5pt] (0.6,4.85) rectangle (12.3,8.53);
    \fill[gray!20] (0.6,8.00) rectangle (12.3,8.53);
  \end{scope}
  \draw[panel] (0.6,4.85) rectangle (12.3,8.53);
  \node[ttl] at (12.1,8.26) {\textarabic{سجل التدقيق}};
  \node[uclbl] at (0.8,8.26) {\EN{UC-06}};

  %---- Toolbar: date range + operation-type dropdown + full-log button ------
  \path[field] (9.55,7.55) rectangle (12.10,7.97);
  \node[lblE] at (11.95,7.76)
    {\textarabic{من: \EN{\_\_/\_\_/\_\_\_\_\_}}};
  \path[field] (6.85,7.55) rectangle (9.40,7.97);
  \node[lblE] at (9.25,7.76)
    {\textarabic{إلى: \EN{\_\_/\_\_/\_\_\_\_\_}}};
  \path[field] (4.00,7.55) rectangle (6.70,7.97);
  \node[lblE] at (6.55,7.76) {\textarabic{نوع العملية: الكل}};
  \node[anchor=west, font=\scriptsize, text=gray!70, inner sep=0pt]
    at (4.18,7.76) {$\blacktriangledown$};
  \node[btn, minimum width=2.2cm, minimum height=0.45cm] at (1.9,7.76)
    {\textarabic{عرض السجل الكامل}};

  %---- Audit table (bidi mirrors tabular: source col1 renders rightmost,
  % so the source is written in the visual right-to-left order: الوقت أولاً) --
  \node[anchor=north east, inner sep=0pt] at (12.1,7.42) {%
    {\footnotesize
    \renewcommand{\arraystretch}{1.3}%
    \begin{tabular}{@{}>{\raggedleft\arraybackslash}p{1.15cm}@{\hspace{6pt}}>{\raggedleft\arraybackslash}p{1.60cm}@{\hspace{6pt}}>{\raggedleft\arraybackslash}p{3.81cm}@{\hspace{6pt}}>{\raggedleft\arraybackslash}p{3.95cm}@{}}
    \arrayrulecolor{gray!50}
    \rowcolor{gray!15}
    \textarabic{\textbf{الوقت}} & \textarabic{\textbf{المستخدم}} &
    \textarabic{\textbf{الإجراء}} & \textarabic{\textbf{النتيجة}} \\ \hline
    \EN{09:15} & \textarabic{أحمد} &
    \textarabic{تسجيل دخول} & \textarabic{نجاح} \\ \hline
    \EN{09:42} & \textarabic{سارة} &
    \textarabic{تحديث دور مستخدم} & \textarabic{نجاح} \\ \hline
    \EN{10:03} & \textarabic{أحمد} &
    \textarabic{تنفيذ استعلام \EN{SQL}} &
    \textcolor{red!70!black}{\textarabic{\textbf{مرفوض (تجاوز الصلاحية)}}} \\ \hline
    \end{tabular}}};

  %==========================================================================
  % PANEL 3 (bottom): العروض الذكية [UC-05] — absent from the original draft
  %==========================================================================
  \path[panel, fill=gray!5, draw=none] (0.6,0.75) rectangle (12.3,4.63);
  \begin{scope}
    \clip[rounded corners=5pt] (0.6,0.75) rectangle (12.3,4.63);
    \fill[gray!20] (0.6,4.10) rectangle (12.3,4.63);
  \end{scope}
  \draw[panel] (0.6,0.75) rectangle (12.3,4.63);
  \node[ttl] at (12.1,4.36) {\textarabic{العروض الذكية}};
  \node[uclbl] at (0.8,4.36) {\EN{UC-05}};

  %---- Toolbar: catalog search + create-view button -------------------------
  \path[field] (7.6,3.53) rectangle (12.1,3.98);
  \node[ph] at (11.95,3.76) {\textarabic{بحث في الكتالوج...}};
  \node[btnprim, minimum width=2.05cm, minimum height=0.45cm]
    at (1.825,3.76) {\textarabic{إنشاء عرض ذكي}};

  %---- Catalog card 1 (rightmost): عام — معتمد ------------------------------
  \path[viewcard] (8.5,0.95) rectangle (12.1,3.4);
  \node[cardtxt] at (10.3,3.10) {\CODE{\textbf{v\_academic\_catalog}}};
  \node[cardtxt] at (10.3,2.74) {\textarabic{دليل البرامج الأكاديمية}};
  \node[cardtxt] at (10.3,2.38) {\textarabic{التصنيف: عام}};
  \node[cardtxt] at (10.3,2.02) {\textarabic{الحالة: معتمد}};
  \node[cbtn] at (11.08,1.40) {\textarabic{مراجعة}};
  \node[cbtn] at (9.52,1.40) {\textarabic{اعتماد}};

  %---- Catalog card 2 (middle): مقنّع — معتمد -------------------------------
  \path[viewcard] (4.65,0.95) rectangle (8.25,3.4);
  \node[cardtxt] at (6.45,3.10) {\CODE{\textbf{v\_student\_profile\_masked}}};
  \node[cardtxt] at (6.45,2.74) {\textarabic{الملف الأساسي المقنّع للطالب}};
  \node[cardtxt] at (6.45,2.38) {\textarabic{التصنيف: مقنّع}};
  \node[cardtxt] at (6.45,2.02) {\textarabic{الحالة: معتمد}};
  \node[cbtn] at (7.23,1.40) {\textarabic{مراجعة}};
  \node[cbtn] at (5.67,1.40) {\textarabic{اعتماد}};

  %---- Catalog card 3 (leftmost): مقيّد — قيد المراجعة (اعتماد معطّل) --------
  \path[viewcard] (0.8,0.95) rectangle (4.4,3.4);
  \node[cardtxt] at (2.6,3.10) {\CODE{\textbf{v\_student\_enrollment}}};
  \node[cardtxt] at (2.6,2.74) {\textarabic{قيد الطالب والبرنامج والدفعة}};
  \node[cardtxt] at (2.6,2.38) {\textarabic{التصنيف: مقيّد}};
  \node[cardtxt] at (2.6,2.02) {\textarabic{الحالة: قيد المراجعة}};
  \node[cbtn] at (3.38,1.40) {\textarabic{مراجعة}};
  \node[cbtn, dashed, fill=white] at (1.82,1.40) {\textarabic{اعتماد}};

  %==========================================================================
  % SCREEN FRAME (stroke on top — closes the wireframe)
  %==========================================================================
  \draw[gray!70, line width=1pt, rounded corners=8pt]
    (0.4,0.5) rectangle (15.6,13.1);

\end{tikzpicture}%
}

  \caption{المخطط السلكي للوحة الإدارة بواجهة عربية من اليمين إلى اليسار، بمهامها الثلاث: إدارة المستخدمين والأدوار (\EN{UC-04})، وسجل التدقيق (\EN{UC-06})، وإدارة العروض الذكية واعتمادها (\EN{UC-05}).}
  \label{fig:wireframes-admin}
\end{figure}

ولم يُرسم في هذين المخططين ما هو دون مستوى الشاشة: فتح قائمة «عرض استعلام \EN{SQL}» وتوسيعها، وحالات التركيز والإدخال، واستجابات الأزرار؛ فهذه تفاصيل تنفيذ تُقاس في فصل الاختبار ضمن اختبار قابلية الاستخدام (\EN{NFR-13}، \EN{NFR-14})، وتُختبر بها صحة الاتجاه من اليمين إلى اليسار (\EN{NFR-15}). وبذلك يكتمل عنصر التصميم الأخير في الفصل، وتُجمع آثاره كلها في مصفوفة التتبع التي تربط كل عنصر تصميم بمتطلبه الأصلي (القسم \ref{sec:design-traceability}).
